\section{设计含义：抽象层次决定可扩展性}
把 Base、Worker、Manager 三层拆开，并不是为了“代码更好看”，而是为了让 Agentic Loop 可以在不同任务间复用。真正变化最快的是具体任务逻辑，而不是状态管理和推理调度，因此统一抽象层能显著降低后续实验成本。

\begin{knowledgebox}{读框架代码时先看什么}
\begin{itemize}
  \item 哪一层负责统一 orchestration
  \item 哪一层承接具体任务状态
  \item 哪些接口是为了让研究者替换业务逻辑
\end{itemize}
\end{knowledgebox}

\subsection{本章小结}
理解 AgentLoopManager 的最佳方式，不是死记函数名，而是看清哪些职责被稳定地抽象出来，哪些职责被故意留给自定义逻辑。

\newpage
